\documentclass[a4paper,12pt]{article}

\usepackage{amsmath,amssymb,amsthm}
\usepackage[margin=25mm]{geometry}
\usepackage{enumitem}
\usepackage{mathtools}

\title{ECAにおけるKの集合論的基礎・数値化・区間評価と$T_{\mathrm{overlay}}$への接続}
\author{}
\date{}

\begin{document}

\maketitle

\section{本稿の目的と位置付け}

本稿では、拡張連続性公理（ECA）における
$K$ の定義と数値化について、Solver集合を基礎とする
集合論的な構成から、区間評価、比較対象間の差、
近似値の採用条件、および
$T_{\mathrm{overlay}}$ への接続までを一貫した形で整理する。

$K$ は、統計学的な分散や統計量として定義するものではない。
その基礎は、Solver集合上における公理の配置、
配置間の関係、およびそこから生じる分離にある。
数値化は、この集合論的な構造を評価するために導入するものであり、
数値化の手続きによって集合論的な基礎を置き換えるものではない。

本稿では、次の段階を区別する。

\begin{enumerate}[label=(\arabic*)]
    \item Solver集合上の公理配置を集合として定義する。
    \item 配置間の関係と分離を集合演算によって記述する。
    \item 公理の配置に対応する数値を定義する。
    \item 最小値・平均値・最大値および方向性のある差を用いて比較する。
    \item 一意な数値を保証できない場合には区間評価を維持する。
    \item 条件が満たされる場合に限り近似値を採用する。
    \item $B$、$I$、$K$ を用いる統合評価
    $T_{\mathrm{overlay}}$ に接続する。
\end{enumerate}

これらの段階のうち、集合論的な定義と区間評価の原則は
本稿の基礎となる。
一方、複数の公理の数値を一つのSolver単位の数値に集約する
一般関数については、特定の方式を最終決定したものとして扱わない。

\section{Solver集合と公理の配置}

\subsection{基本となる集合}

$\mathcal{A}$ を公理の集合、$\mathcal{Z}$ をSolver集合とする。
各 $z\in\mathcal{Z}$ に対して、$z$ に対応する公理集合を
$\mathcal{A}_z\subseteq\mathcal{A}$ とする。

公理 $a\in\mathcal{A}$ が含まれるSolverの集合を
\[
    \mathcal{Z}_a
    =
    \{z\in\mathcal{Z}\mid a\in\mathcal{A}_z\}
\]
と定義する。

したがって、$\mathcal{Z}_a$ は、公理 $a$ が
どのSolverに配置されているかを表す集合である。
$K$ の集合論的な検討では、この配置集合を出発点とする。

複数の公理からなる集合 $A\subseteq\mathcal{A}$ に対しては、
\[
    \mathcal{Z}_A
    =
    \bigcap_{a\in A}\mathcal{Z}_a
    =
    \{z\in\mathcal{Z}\mid A\subseteq\mathcal{A}_z\}
\]
と定義する。
すなわち、$\mathcal{Z}_A$ は、$A$ に含まれるすべての公理を
含むSolverの集合である。

なお、$A=\varnothing$ の場合には空集合上の共通部分を
$\mathcal{Z}$ とする。
ただし、複数の公理を比較する際には、比較対象となる公理集合が
空でないことを別途確認する。

\subsection{公理からSolver集合の部分集合への写像}

公理の配置を写像として記述するため、
\[
    \Phi_{\mathcal{Z}}:
    \mathcal{A}\longrightarrow\mathcal{P}(\mathcal{Z})
\]
を
\[
    \Phi_{\mathcal{Z}}(a)=\mathcal{Z}_a
\]
によって定義する。
ここで、$\mathcal{P}(\mathcal{Z})$ は
$\mathcal{Z}$ のべき集合である。

この写像は、公理をその配置集合に対応付けるものである。
$\Phi_{\mathcal{Z}}$ 自体を $K$ と同一視するのではなく、
$K$ を構成するための集合論的な基盤として扱う。

\section{配置・関係・分離の集合論的記述}

\subsection{共通配置と差集合}

公理 $a,b\in\mathcal{A}$ に対して、
両者の配置関係を次の集合によって記述する。
\[
    C(a,b)=\mathcal{Z}_a\cap\mathcal{Z}_b,
\]
\[
    D(a\mid b)=\mathcal{Z}_a\setminus\mathcal{Z}_b,
\]
\[
    D(b\mid a)=\mathcal{Z}_b\setminus\mathcal{Z}_a.
\]

$C(a,b)$ は両方の公理を含むSolverの集合である。
$D(a\mid b)$ は $a$ を含むが $b$ を含まないSolverの集合を表し、
$D(b\mid a)$ はその逆を表す。

したがって、これらの集合は単なる公理間の数値差ではなく、
Solver集合上における共通配置と方向別の分離を記述する。

\subsection{対称差による分離}

両公理の配置の違いを方向に依存しない形で記述するため、
対称差
\[
    S(a,b)
    =
    \mathcal{Z}_a\triangle\mathcal{Z}_b
\]
を定義する。ここで、
\[
    \mathcal{Z}_a\triangle\mathcal{Z}_b
    =
    D(a\mid b)\cup D(b\mid a)
\]
である。

$S(a,b)$ は、一方の公理だけを含むSolverをすべて集めた集合である。
したがって、$S(a,b)=\varnothing$ であることと
$\mathcal{Z}_a=\mathcal{Z}_b$ であることは同値である。

一方、
\[
    C(a,b)=\varnothing
\]
である場合、両公理を同時に含むSolverは存在しない。
これは共通配置が存在しないことを表すが、
それだけで他のすべての配置関係が決まるわけではない。

\subsection{包含関係}

配置集合の包含関係も重要である。例えば、
\[
    \mathcal{Z}_a\subseteq\mathcal{Z}_b
\]
であれば、$a$ を含むすべてのSolverが $b$ も含む。
これはSolver集合上の配置に関する含意を表す。

ただし、配置集合の包含関係から、公理そのものの論理的な同値性や
一般的な証明可能性を無条件に導くことはできない。
そのような結論には、別途、論理体系および公理の意味に関する条件が必要である。

\subsection{$K$ の基本的な解釈}

以上から、$K$ が記述すべき内容は、次の三つを統合したものである。

\begin{enumerate}[label=(\arabic*)]
    \item \textbf{配置：}
    各公理がSolver集合のどの部分に対応するか。
    \item \textbf{関係：}
    複数の公理の配置集合がどの程度共通し、
    どのような包含関係を持つか。
    \item \textbf{分離：}
    配置集合の差集合および対称差がどのように構成されるか。
\end{enumerate}

この三つは独立した代替定義ではなく、
Solver集合上の公理配置を多面的に記述するための相互に関連した要素である。

\section{公理配置の数値化}

\subsection{公理ごとの配置値}

集合論的な配置を数値として評価するため、
まず各公理 $a\in\mathcal{A}$ に対する配置値 $K(a)$ を考える。

$\mathcal{Z}$ が有限かつ空でない場合には、
\[
    K(a)=\frac{|\mathcal{Z}_a|}{|\mathcal{Z}|}
\]
と定義する。
ここで、$|\cdot|$ は集合の要素数を表す。
この定義では $K(a)\in[0,1]$ であり、
公理 $a$ がSolver集合全体のうちどの割合のSolverに配置されているかを表す。

この定義から、
\[
    K(a)=0
    \quad\Longleftrightarrow\quad
    \mathcal{Z}_a=\varnothing
\]
および
\[
    K(a)=1
    \quad\Longleftrightarrow\quad
    \mathcal{Z}_a=\mathcal{Z}
\]
が成り立つ。

したがって、有限なSolver集合に対するこの数値化では、
$K(a)=0$ は公理 $a$ を含むSolverが存在しないことを、
$K(a)=1$ はすべてのSolverが公理 $a$ を含むことを表す。

\subsection{測度を用いる場合}

$\mathcal{Z}$ が有限集合でない場合などには、
Solver集合上の測度空間
\[
    (\mathcal{Z},\Sigma_{\mathcal{Z}},\mu)
\]
を考えることができる。
ここで、$\Sigma_{\mathcal{Z}}$ は可測集合からなる
$\sigma$-代数、$\mu$ はその上で定義された測度である。

$\mathcal{Z}_a$ が可測であり、
\[
    0<\mu(\mathcal{Z})<\infty
\]
が成り立つ場合には、
\[
    K(a)=\frac{\mu(\mathcal{Z}_a)}{\mu(\mathcal{Z})}
\]
と定義できる。この場合も $K(a)\in[0,1]$ である。

この式を使用するには、配置集合の可測性と分母の条件が必要である。
$\sigma$-代数が定められていることだけでは、
$\mathcal{Z}$ のすべての部分集合が可測であるとは限らない。

また、測度による数値化では、測度がゼロである集合と空集合を
同一視してはならない。例えば、$\mu(\mathcal{Z}_a)=0$ であっても、
$\mathcal{Z}_a=\varnothing$ とは限らない。

\subsection{複数の公理を含むSolverの評価}

各Solver $z\in\mathcal{Z}$ に含まれる公理集合を
$\mathcal{A}_z$ とする。
公理ごとの配置値からSolver単位の値を構成する場合、
\[
    K(z)=F\bigl(\{K(a)\mid a\in\mathcal{A}_z\}\bigr)
\]
という形式が考えられる。

ここで $F$ は、複数の公理の配置値をSolver単位の評価へ集約する関数である。
ただし、$F$ を平均・最小・最大のいずれか一つに固定することは、
本稿で確定した内容には含めない。

最小値・平均値・最大値は、Solver内の公理配置を比較するための
具体的な指標として定義できる。
これらを最終的な $K(z)$ にどのように集約するかは、
それぞれの指標が表す情報を踏まえて別途判断する必要がある。

\section{Solver内の最小値・平均値・最大値}

\subsection{基本定義}

$\mathcal{A}_z$ が有限かつ空でないとし、
\[
    \mathcal{A}_z=\{a_1,\ldots,a_n\}
\]
とする。このとき、Solver $z$ に含まれる公理の配置値について、
\[
    K_{\min}(z)
    =
    \min_{a\in\mathcal{A}_z}K(a),
\]
\[
    K_{\mathrm{mean}}(z)
    =
    \frac{1}{|\mathcal{A}_z|}
    \sum_{a\in\mathcal{A}_z}K(a),
\]
\[
    K_{\max}(z)
    =
    \max_{a\in\mathcal{A}_z}K(a)
\]
と定義する。

これらはそれぞれ、Solverに含まれる公理の配置値の最小値、
平均値、最大値を表す。
通常の有限集合上の実数値として、
\[
    K_{\min}(z)
    \leq K_{\mathrm{mean}}(z)
    \leq K_{\max}(z)
\]
が成り立つ。

\subsection{配置値の範囲と平均からの差}

Solver内の配置値の広がりを記述するため、
\[
    R_K(z)=K_{\max}(z)-K_{\min}(z)
\]
と定義する。

さらに、平均値と最小値の差、および最大値と平均値の差を
\[
    \delta_-(z)
    =
    K_{\mathrm{mean}}(z)-K_{\min}(z),
\]
\[
    \delta_+(z)
    =
    K_{\max}(z)-K_{\mathrm{mean}}(z)
\]
と定義する。

これらには、
\[
    R_K(z)=\delta_-(z)+\delta_+(z)
\]
という関係がある。

ただし、$R_K(z)$、$\delta_-(z)$、$\delta_+(z)$ は
異なる特徴を記述する。
範囲 $R_K(z)$ は最大値と最小値の隔たりを表し、
$\delta_-(z)$ と $\delta_+(z)$ は平均値が両端から
どの程度離れているかを表す。
したがって、これらの値を一つの意味に還元せず、
相互に関連する指標として扱う。

\section{公理間の数値比較}

\subsection{絶対差と方向性のある差}

公理 $a,b\in\mathcal{A}$ の配置値が確定している場合、
両者の絶対差を
\[
    d(a,b)=|K(a)-K(b)|
\]
と定義する。

また、方向性のある差を
\[
    d_{\to}(a,b)=K(a)-K(b)
\]
と定義する。このとき、
\[
    d_{\to}(a,b)=-d_{\to}(b,a)
\]
が成り立つ。

絶対差は比較の方向に依存しない隔たりを表す一方、
方向性のある差は、どちらの配置値が大きいかを保持する。
両者は異なる情報を持つため、必要に応じて併用する。

\subsection{比較対象集合における最小差・平均差・最大差}

有限な比較対象集合 $A_0\subseteq\mathcal{A}$ を考える。
$|A_0|\geq2$ とし、$a\ne b$ のすべての組を比較対象とする。

最小差を
\[
    d_{\min}(A_0)
    =
    \min_{\substack{a,b\in A_0\\a\ne b}}d(a,b)
\]
と定義する。

順序を持たない公理対を一度ずつ数える集合を
\[
    P(A_0)
    =
    \bigl\{\{a,b\}\subseteq A_0\mid a\ne b\bigr\}
\]
とすると、平均差は
\[
    d_{\mathrm{mean}}(A_0)
    =
    \frac{1}{|P(A_0)|}
    \sum_{\{a,b\}\in P(A_0)}d(a,b)
\]
と定義できる。

最大差は
\[
    d_{\max}(A_0)
    =
    \max_{\substack{a,b\in A_0\\a\ne b}}d(a,b)
\]
と定義する。

したがって、
\[
    d_{\min}(A_0)
    \leq d_{\mathrm{mean}}(A_0)
    \leq d_{\max}(A_0)
\]
が成り立つ。

ここで、最小差・平均差・最大差は、単に同じ数値を異なる形で
表示するものではない。
最小差は最も近い比較対象間の隔たり、
平均差は比較対象全体にわたる平均的な隔たり、
最大差は最も離れた比較対象間の隔たりをそれぞれ記述する。

\subsection{比較対象ごとの方向性}

比較対象 $a\in A_0$ を固定した場合、
他の公理との方向性のある差を
\[
    d_{\to}(a,b)=K(a)-K(b),
    \qquad b\in A_0\setminus\{a\}
\]
によって調べることができる。

比較対象 $a$ から見た最小絶対差・平均絶対差・最大絶対差は、
\[
    d_{\min}(a;A_0)
    =
    \min_{b\in A_0\setminus\{a\}}|K(a)-K(b)|,
\]
\[
    d_{\mathrm{mean}}(a;A_0)
    =
    \frac{1}{|A_0|-1}
    \sum_{b\in A_0\setminus\{a\}}|K(a)-K(b)|,
\]
\[
    d_{\max}(a;A_0)
    =
    \max_{b\in A_0\setminus\{a\}}|K(a)-K(b)|
\]
と定義できる。

これらの指標を比較する際には、
比較対象全体についての集約値と、
特定の公理から見た比較値を区別する必要がある。

\section{可測性と区間評価}

\subsection{可測性は測度とは異なる}

Solver集合上の配置集合を数値化する場合、
集合が可測であることと、その集合の測度が定義されていることを
区別する必要がある。

測度空間 $(\mathcal{Z},\Sigma_{\mathcal{Z}},\mu)$ において、
$\Sigma_{\mathcal{Z}}$ に属さない集合に対して、
$\mu$ をそのまま通常の測度として適用することはできない。

したがって、$\mathcal{Z}_a$ や
$\mathcal{Z}_a\triangle\mathcal{Z}_b$ が非可測である場合には、
可測集合に対する測度の式を無条件に適用してはならない。

また、ある集合が非可測であるかどうかは、
採用した $\sigma$-代数および測度空間に依存する。
非可測性を、測度構造と無関係な絶対的性質として扱わない。

\subsection{外測度と内測度}

一般の部分集合 $E\subseteq\mathcal{Z}$ を評価するため、
外測度と内測度を
\[
    \mu^*(E)
    =
    \inf\{\mu(M)\mid M\in\Sigma_{\mathcal{Z}},\
    E\subseteq M\},
\]
\[
    \mu_*(E)
    =
    \sup\{\mu(N)\mid N\in\Sigma_{\mathcal{Z}},\
    N\subseteq E\}
\]
とする。

これらの定義では、外測度は $E$ を包含する可測集合の測度の
下限を、内測度は $E$ に含まれる可測集合の測度の上限を表す。

$\mu(\mathcal{Z})=1$ である確率測度など、
全体の測度が $1$ の有限測度空間を考えると、
\[
    0\leq\mu_*(E)\leq\mu^*(E)\leq1
\]
が成り立つ。

このとき、
\[
    \mu_*(E)\leq\mu(E)\leq\mu^*(E)
\]
という通常の区間評価は、$E$ が可測である場合に限って
通常の測度 $\mu(E)$ と結び付く。
非可測集合に対しては、$\mu(E)$ が通常の測度として
定義されているかのように記述してはならない。

\subsection{外測度と内測度の差}

区間幅に相当する量として、
\[
    \Delta_{\mu}(E)
    =
    \mu^*(E)-\mu_*(E)
\]
を定義する。

$\Delta_{\mu}(E)$ は、外測度と内測度の間に残る幅を表す。
$E$ が可測であれば両者は一致し、
\[
    \Delta_{\mu}(E)=0
\]
となる。

一方、一般の測度空間では、$\Delta_{\mu}(E)=0$ であることは、
$E$ がその測度空間の完備化において可測であることに対応する。
元の $\sigma$-代数に関する可測性と、完備化における可測性は
区別しなければならない。

非可測集合に対しては、外測度と内測度が一致しないことがあり、
その場合には区間の幅が正になる。
この幅をゼロと仮定したり、任意に小さくできると仮定したりしてはならない。

\section{区間値としての公理配置評価}

\subsection{公理ごとの区間評価}

公理 $a$ に対する配置値が一意に定まらない場合、
単一の実数値を無理に割り当てず、
\[
    K(a)\in[k_a^-,k_a^+]
\]
と評価する。

区間幅を
\[
    \Delta K(a)=k_a^+-k_a^-
\]
と定義する。

$k_a^-$ と $k_a^+$ は、それぞれ配置値の下限と上限である。
区間評価は、値が未確定であることを隠すためのものではなく、
現時点で保証できる範囲を明示するためのものである。

例えば、正規化された有限測度空間で、
配置集合に対して外測度・内測度による評価を採用する場合には、
\[
    k_a^-=\mu_*(\mathcal{Z}_a),
    \qquad
    k_a^+=\mu^*(\mathcal{Z}_a)
\]
とすることができる。
この式は、$\mathcal{Z}$ 全体の測度が $1$ である場合のものであり、
正規化されていない有限測度空間では、
全体の測度による正規化が別途必要になる。

\subsection{Solver内の最小値・平均値・最大値の区間評価}

$\mathcal{A}_z=\{a_1,\ldots,a_n\}$ が有限かつ空でなく、
各公理の配置値が
\[
    K(a_i)\in[k_i^-,k_i^+]
\]
と評価されているとする。

最小値の区間評価は
\[
    K_{\min}(z)
    \in
    \left[
        \min_i k_i^-,
        \min_i k_i^+
    \right]
\]
である。

平均値の区間評価は
\[
    K_{\mathrm{mean}}(z)
    \in
    \left[
        \frac{1}{n}\sum_{i=1}^n k_i^-,
        \frac{1}{n}\sum_{i=1}^n k_i^+
    \right]
\]
である。

最大値の区間評価は
\[
    K_{\max}(z)
    \in
    \left[
        \max_i k_i^-,
        \max_i k_i^+
    \right]
\]
である。

これらは各配置値の区間だけから得られる安全な包含区間である。
ただし、複数の区間値の間に依存関係がある場合には、
区間端点のすべてが同時に実現可能とは限らない。
そのため、得られた区間は実際に取り得る値の範囲より
広くなる場合がある。

\subsection{範囲と平均からの差の区間評価}

\[
    R_K(z)=K_{\max}(z)-K_{\min}(z)
\]
については、
\[
    R_K(z)\in
    \left[
        \max\left\{0,\,
        \max_i k_i^--\min_i k_i^+\right\},
        \max_i k_i^+-\min_i k_i^-
    \right]
\]
と評価できる。

同様に、
\[
    \delta_-(z)=K_{\mathrm{mean}}(z)-K_{\min}(z)
\]
について、
\[
    \delta_-(z)\in
    \left[
        \max\{0,K_{\mathrm{mean}}^--K_{\min}^+\},
        K_{\mathrm{mean}}^+-K_{\min}^-
    \right]
\]
とする。

また、
\[
    \delta_+(z)=K_{\max}(z)-K_{\mathrm{mean}}(z)
\]
について、
\[
    \delta_+(z)\in
    \left[
        \max\{0,K_{\max}^--K_{\mathrm{mean}}^+\},
        K_{\max}^+-K_{\mathrm{mean}}^-
    \right]
\]
とする。

ここで、$K_{\min}^-$、$K_{\min}^+$、
$K_{\mathrm{mean}}^-$、$K_{\mathrm{mean}}^+$、
$K_{\max}^-$、$K_{\max}^+$ は、
それぞれの区間評価の下限と上限を表す。

これらの式は、区間演算による保守的な評価である。
依存関係を考慮しない場合には、区間の端点が同時に実現するとは限らない。
したがって、区間演算によって得られた端点を、
実際の最小値・最大値として無条件に解釈しない。

\section{区間評価における公理間の差}

\subsection{方向性のある差の区間}

\[
    K(a)\in[k_a^-,k_a^+],
    \qquad
    K(b)\in[k_b^-,k_b^+]
\]
であるとする。

方向性のある差は、
\[
    d_{\to}(a,b)=K(a)-K(b)
\]
であるから、
\[
    d_{\to}(a,b)
    \in
    [k_a^--k_b^+,\,k_a^+-k_b^-]
\]
となる。

この区間が正の値だけを含む場合、
$K(a)>K(b)$ を保証できる。
負の値だけを含む場合には $K(a)<K(b)$ を保証できる。
一方、区間が $0$ を含む場合には、
この区間評価だけから大小関係を確定できない。

\subsection{絶対差の区間}

絶対差
\[
    d(a,b)=|K(a)-K(b)|
\]
について、下限と上限を
\[
    \ell_{ab}
    =
    \max\{0,\ k_a^--k_b^+,\ k_b^--k_a^+\},
\]
\[
    u_{ab}
    =
    \max\{
        |k_a^--k_b^+|,
        |k_a^+-k_b^-|
    \}
\]
と定義する。

このとき、
\[
    d(a,b)\in[\ell_{ab},u_{ab}]
\]
である。

$\ell_{ab}>0$ ならば、両者の差が正であることが保証される。
一方、$\ell_{ab}=0$ である場合には、
区間評価から正の差を保証できない。

\subsection{比較対象全体の最小差の区間}

有限な比較対象集合 $A_0$ に対し、
$|A_0|\geq2$ とする。
各公理対の絶対差が
\[
    d(a,b)\in[\ell_{ab},u_{ab}]
\]
である場合、最小差は
\[
    d_{\min}(A_0)
    \in
    \left[
        \min_{\substack{a,b\in A_0\\a\ne b}}\ell_{ab},
        \min_{\substack{a,b\in A_0\\a\ne b}}u_{ab}
    \right]
\]
と評価できる。

ここで、下限と上限を
\[
    d_{\min}^-(A_0)
    =
    \min_{\substack{a,b\in A_0\\a\ne b}}\ell_{ab},
\]
\[
    d_{\min}^+(A_0)
    =
    \min_{\substack{a,b\in A_0\\a\ne b}}u_{ab}
\]
と書く。

したがって、
\[
    d_{\min}(A_0)
    \in[d_{\min}^-(A_0),d_{\min}^+(A_0)]
\]
である。

\section{差の下限がゼロであることの意味}

\subsection{三つの状態の区別}

区間評価によって差の下限がゼロとなる場合、
それを直ちに「実際の差がゼロである」と解釈してはならない。
少なくとも、次の三つを区別する必要がある。

\begin{enumerate}[label=(\arabic*)]
    \item \textbf{差がゼロであると確定している状態：}
    $K(a)=K(b)$ が確定している。
    \item \textbf{差がゼロである可能性を排除できない状態：}
    区間が重なっており、$K(a)=K(b)$ の可能性が残る。
    \item \textbf{差が正であると保証されている状態：}
    $\ell_{ab}>0$ であり、両者の差が正であることが保証される。
\end{enumerate}

第1の場合と第2の場合は異なる。
区間が一点に定まって両者が同じ値を取る場合には、
差がゼロであることが確定する。
一方、幅のある区間が重なっている場合には、
実際の値が異なっている可能性も残る。

\subsection{最小差の正値保証不能}

比較対象集合 $A_0$ に対して、
\[
    d_{\min}^-(A_0)=0
\]
である場合を、\textbf{最小差の正値保証不能}と呼ぶ。

これは、比較対象間の最小差について、
区間評価から正の下限を保証できない状態を意味する。
実際の最小差がゼロであることを意味するものではない。

また、$d_{\min}^-(A_0)=0$ であることから、
比較対象のすべての組が区別不能であると結論してはならない。
一部の組については正の差が保証されていても、
別の組について正の差が保証されない可能性がある。

したがって、区間評価の結果はそのまま維持し、
正の差が保証されていないことを明示する。
これによって、値が一致しているという結論と、
一致の可能性を排除できないという結論を区別できる。

\section{許容誤差と近似値の採用条件}

\subsection{区間評価を基本とする原則}

公理 $a$ の配置値が
\[
    K(a)\in[k_a^-,k_a^+]
\]
で与えられる場合、区間幅
\[
    \Delta K(a)=k_a^+-k_a^-
\]
を用いて評価の幅を表す。

基本方針は、区間評価を標準とすることである。
区間の幅が正であるという理由だけで、
一つの実数値を採用することはしない。

単一の近似値が必要となる場合には、
あらかじめ定めた許容誤差 $\varepsilon>0$ を用いて、
区間幅が条件を満たすかどうかを確認する。

\subsection{中点による近似}

区間
\[
    K(a)\in[k_a^-,k_a^+]
\]
に対する中点近似を
\[
    \widehat{K}(a)
    =
    \frac{k_a^-+k_a^+}{2}
\]
と定義する。

この近似値について、
\[
    |K(a)-\widehat{K}(a)|
    \leq\frac{\Delta K(a)}{2}
\]
が成り立つ。

したがって、区間幅が十分小さく、単一の数値が必要である場合には、
中点を近似値として用いることができる。
ただし、中点を採用することは、
区間内の真の値が中点に等しいと主張することではない。

\subsection{最小差による区別可能性の確認}

有限な比較対象集合 $A_0$ に対し、
\[
    \Delta_{\max}
    =
    \max_{a\in A_0}\Delta K(a)
\]
と定義する。

また、比較対象間の差の下限に基づき、
\[
    d_{\min}^-(A_0)
    =
    \min_{\substack{a,b\in A_0\\a\ne b}}\ell_{ab}
\]
とする。

このとき、
\[
    \Delta_{\max}<d_{\min}^-(A_0)
\]
が成り立つならば、各公理の中点近似に伴う誤差は、
比較対象間の差の保証された下限より小さい。

実際、
\[
    |\widehat{K}(a)-K(a)|
    \leq\frac{\Delta_{\max}}{2}
\]
であり、任意の異なる $a,b\in A_0$ について、
\[
\begin{aligned}
    |\widehat{K}(a)-\widehat{K}(b)|
    &\geq |K(a)-K(b)|
       -|\widehat{K}(a)-K(a)|\\
    &\qquad
       -|\widehat{K}(b)-K(b)|\\
    &\geq \ell_{ab}-\Delta_{\max}\\
    &\geq d_{\min}^-(A_0)-\Delta_{\max}\\
    &>0
\end{aligned}
\]
となる。

したがって、この条件の下では、
中点近似によって比較対象間の区別が失われることはない。

ただし、この条件は十分条件であり、必要条件ではない。
条件を満たさない場合でも、特定の比較対象間の区別が保たれることはあり得る。

\subsection{最小差の下限がゼロの場合}

\[
    d_{\min}^-(A_0)=0
\]
の場合、正の区間幅を持つ評価について、
上記の条件
\[
    \Delta_{\max}<d_{\min}^-(A_0)
\]
を満たすことはできない。

この場合には、区間評価を維持する。
中点などの近似値を計算すること自体は可能だが、
比較対象間の区別が保証されているとはいえない。

この状態を、実際の差がゼロであることと混同せず、
「最小差の正値保証不能」として記録する。

\subsection{許容誤差による採用手順}

近似値を採用する場合には、少なくとも次の条件を区別して扱う。

\begin{enumerate}[label=(\arabic*)]
    \item 許容誤差 $\varepsilon>0$ を事前に定める。
    \item 対象となる区間幅が許容誤差の条件を満たすか確認する。
    \item 比較対象がある場合には、最小差の下限を確認する。
    \item 最小差の正値保証が得られない場合には、
    区別可能性を保証できない旨を明示する。
    \item 条件を満たして単一値が必要な場合には、
    中点近似を用いることができる。
    \item 条件を満たさない場合には、区間評価を維持する。
\end{enumerate}

最小差・平均差・最大差、および方向性のある差は、
近似によってどのような比較情報が失われ得るかを確認するために用いる。
単一の許容誤差条件だけで、これらすべての情報が自動的に保存されると
仮定してはならない。

\section{$B$、$I$、$K$ による統合評価への接続}

\subsection{離散形式}

ECAの統合評価を
\[
    T_{\mathrm{overlay}}
    =
    \frac{1}{N}\sum_{i=1}^{N}B_iI_iK_i
\]
とする。

ここで、$B_i$、$I_i$、$K_i$ は各評価対象に対応する値である。
$K_i$ は、本稿で整理した公理配置の評価に対応する。
この式では、$B$、$I$、$K$ はそれぞれ異なる役割を持ち、
$K$ の区間評価を導入したからといって、
$B$ や $I$ の意味を変更するものではない。

\subsection{$K_i$ が区間値である場合}

各項の $B_i$ と $I_i$ が確定した非負の実数であり、
\[
    K_i\in[K_i^-,K_i^+]
\]
であるとする。

この場合、
\[
    T_{\mathrm{overlay}}
    \in[T^-,T^+]
\]
であり、
\[
    T^-=
    \frac{1}{N}\sum_{i=1}^{N}B_iI_iK_i^-,
\]
\[
    T^+=
    \frac{1}{N}\sum_{i=1}^{N}B_iI_iK_i^+
\]
と定義できる。

区間幅は
\[
\begin{aligned}
    \Delta T
    &=T^+-T^-\\
    &=\frac{1}{N}
      \sum_{i=1}^{N}B_iI_i(K_i^+-K_i^-)\\
    &=\frac{1}{N}
      \sum_{i=1}^{N}B_iI_i\Delta K_i
\end{aligned}
\]
となる。

したがって、各 $K_i$ の区間幅は、
$B_iI_i$ による重み付けを通じて
$T_{\mathrm{overlay}}$ の区間幅へ伝播する。

$B_iI_i$ が大きい項ほど、その項の $K_i$ の区間幅が
統合評価に与える寄与は大きい。
一方、$B_iI_i=0$ である項は、
この式において $K_i$ の区間幅を
$T_{\mathrm{overlay}}$ の幅へ寄与させない。

この評価は、$B_iI_i\geq0$ が成立することに依存する。
$B_i$ や $I_i$ が区間値である場合、
または負の値を取り得る場合には、
積の下限と上限を別途評価しなければならない。

\subsection{統合評価の区間幅と近似値}

$T_{\mathrm{overlay}}$ が
\[
    T_{\mathrm{overlay}}\in[T^-,T^+]
\]
と評価された場合、その区間幅を
\[
    \Delta T=T^+-T^-
\]
とする。

単一の実数値が必要で、かつ許容誤差の条件を満たす場合には、
\[
    \widehat{T}_{\mathrm{overlay}}
    =
    \frac{T^-+T^+}{2}
\]
を近似値として採用できる。

このとき、
\[
    |T_{\mathrm{overlay}}-\widehat{T}_{\mathrm{overlay}}|
    \leq\frac{\Delta T}{2}
\]
が成り立つ。

ただし、$K$ の段階で用いた許容誤差を、
何の検討もなく $T_{\mathrm{overlay}}$ にそのまま適用するとは限らない。
統合評価の近似値に必要な精度は、
統合評価そのものの比較目的および許容誤差に基づいて判断する。

\subsection{統合評価間の比較とゼロの解釈}

複数の統合評価値を区間として比較する場合にも、
差の下限がゼロであることを、
実際の値が一致していることと混同してはならない。

例えば、
\[
    T_1\in[T_1^-,T_1^+],
    \qquad
    T_2\in[T_2^-,T_2^+]
\]
の場合、方向性のある差は
\[
    T_1-T_2
    \in
    [T_1^--T_2^+,\,T_1^+-T_2^-]
\]
と評価できる。

この区間がゼロを含む場合、
区間評価だけから $T_1>T_2$ または $T_1<T_2$ を
確定することはできない。

絶対差の下限がゼロである場合も同様に、
差がゼロである可能性を排除できないことを意味するのであって、
実際の差がゼロであることを確定するものではない。

\section{連続形式への接続}

\subsection{連続形式の基本式}

連続形式の統合評価を
\[
    T_{\mathrm{overlay}}
    =
    \int_0^1 B(x)I(x)K(x)\,dx
\]
とする。

この式を通常の積分として用いるには、
被積分関数の可測性および積分可能性などの条件が必要である。
測度を明示する場合には、
\[
    T_{\mathrm{overlay}}
    =
    \int_{[0,1]}B(x)I(x)K(x)\,d\mu(x)
\]
と表すことができる。

ここで、離散形式と連続形式は
統合評価の異なる評価形式である。
K の集合論的な基礎や、区間評価に関する解釈原則を
変更するものではない。

\subsection{連続形式での区間評価}

可測な関数 $K^-(x)$、$K^+(x)$ によって
\[
    K^-(x)\leq K(x)\leq K^+(x)
\]
が成り立つとする。

さらに、$B(x)I(x)\geq0$ がほとんど至る所で成立し、
対応する関数が積分可能であるとき、
\[
    T_{\mathrm{overlay}}
    \in[T^-,T^+]
\]
と評価できる。ただし、
\[
    T^-=
    \int_0^1B(x)I(x)K^-(x)\,dx,
\]
\[
    T^+=
    \int_0^1B(x)I(x)K^+(x)\,dx
\]
である。

このとき、
\[
\begin{aligned}
    \Delta T
    &=T^+-T^-\\
    &=\int_0^1
      B(x)I(x)\bigl(K^+(x)-K^-(x)\bigr)\,dx
\end{aligned}
\]
となる。

ただし、各点で区間端点が定義されていることだけでは、
これらの関数の可測性や積分可能性は保証されない。
その条件が確認できない場合には、
通常の積分によって区間端点を評価できると無条件に主張してはならない。

また、$K$ の評価に外測度・内測度を用いる場合、
それらの点ごとの評価が直ちに積分可能な上下界を与えるとは限らない。
連続形式への接続では、積分に必要な条件を別途確認する。

\section{本稿における確定事項と未確定事項}

\subsection{確定した基本方針}

本稿で整理した基本方針は、次のとおりである。

\begin{enumerate}[label=(\arabic*)]
    \item $K$ はSolver集合上の公理配置・関係・分離を記述する。
    \item 公理 $a$ の配置集合は
    $\mathcal{Z}_a=\{z\in\mathcal{Z}\mid a\in\mathcal{A}_z\}$
    とする。
    \item 複数の公理の共通配置、方向別の差、および対称差を
    集合演算によって記述する。
    \item 有限なSolver集合では、配置値の基本形として
    $K(a)=|\mathcal{Z}_a|/|\mathcal{Z}|$ を用いる。
    \item 適切な測度条件の下では、
    $K(a)=\mu(\mathcal{Z}_a)/\mu(\mathcal{Z})$ を用いる。
    \item 配置値を一意に保証できない場合には区間評価を維持する。
    \item 区間幅が正であることと、実際の値が一意でないことを
    同一視しない。
    \item 差の下限がゼロであることは、
    正の差を保証できないことを表す。
    \item 近似値を採用する場合には、許容誤差と比較対象間の最小差を
    確認し、区別可能性を保証できない場合にはその旨を明示する。
    \item $K$ の区間評価は、条件を確認したうえで
    $T_{\mathrm{overlay}}$ の区間評価に伝播させる。
\end{enumerate}

\subsection{本稿で一意に決定していない事項}

一方、以下は本稿によって一意に決定したものではない。

\begin{enumerate}[label=(\arabic*)]
    \item 複数の公理の配置値から、最終的な $K(z)$ を構成する
    集約関数 $F$ の具体的な形式。
    \item 配置・関係・分離の各情報を最終的な $K$ の数値に
    どのように統合するかという具体的な集約式。
    \item あらゆる対象に共通して適用する許容誤差 $\varepsilon$ の
    一意な数値。
    \item すべての比較目的に共通する近似値の採用基準。
\end{enumerate}

これらを未確定のまま区別しておくことは、
本稿で確定した集合論的な基礎や区間評価の原則を
曖昧にするものではない。
むしろ、既に確定した定義と、追加の条件や選択を必要とする部分を
明確に分離するために必要である。

\section{Kの役割と数値化における補足}

\subsection{Kが担う五つの役割}

Kは、Solver集合上の公理配置・関係・分離を記述する指標として、
少なくとも次の五つの役割を持つ。

\begin{enumerate}[label=(\arabic*)]
    \item \textbf{公理配置の傾向の記述：}
    例えば、公理 $a,b$ の配置が互いに近く、
    公理 $c$ の配置がそれらから離れているといった傾向を記述する。

    \item \textbf{対象公理の配置位置の評価：}
    対象となる公理が、比較対象となる公理や配置全体に対して
    どのような位置にあるかを評価する。

    \item \textbf{複合Solver集合における指標：}
    複数の公理体系や配置が関係する複合Solver集合
    $\mathcal{Z}$ において、公理配置を評価するための指標として用いる。

    \item \textbf{Iとの組合せによる要素の抽出：}
    $I$ と組み合わせることで、複合Solver集合から
    条件に対応する要素を抽出するために用いる。
    この役割において、$B$ が常に必要であるとは限らない。

    \item \textbf{Bとの組合せによる除外範囲の評価：}
    Kによる分離が明確に離散的に扱える場合には、
    $B$ と組み合わせることで、$I$ に基づく評価から
    どの範囲を除外できるかを見積もる。
\end{enumerate}

これらはKの用途を列挙したものであり、
五つの役割を単純に加算することによってKの値を定義するものではない。
それぞれの役割がどのような数学的条件の下で成立するかを区別し、
最終的な数値化に接続する必要がある。

\subsection{Kの対象とする公理集合}

Kが対象とするのは、各Solver $z$ に対応する公理集合
$\mathcal{A}_z$ である。
したがって、本稿で扱う配置集合は
\[
    \mathcal{Z}_a
    =
    \{z\in\mathcal{Z}\mid a\in\mathcal{A}_z\}
\]
であり、Kを推論規則や構成制約の集合
$\mathcal{R}_z$ に直接適用することを前提としない。

$\mathcal{R}_z$ が別の箇所で推論規則・構成制約を表す場合、
$\mathcal{A}_z$ と $\mathcal{R}_z$ は役割の異なる対象として区別する。
両者の間に関係を設定する場合には、その関係を別途定義する必要がある。

\subsection{配置値と配置・関係・分離の統合}

公理ごとの配置値
\[
    K(a)=\frac{|\mathcal{Z}_a|}{|\mathcal{Z}|}
\]
または適切な測度条件の下での
\[
    K(a)=\frac{\mu(\mathcal{Z}_a)}{\mu(\mathcal{Z})}
\]
は、公理がSolver集合全体のどの割合に配置されているかを表す。

ただし、これらの式だけでは、配置集合間の関係や分離のすべてが
数値として表現されているわけではない。
例えば、配置値が等しい二つの公理について、
\[
    K(a)=K(b)
\]
が成立していても、
\[
    \mathcal{Z}_a=\mathcal{Z}_b
\]
が必ず成立するわけではない。
異なる配置集合が同じ測度を持つ場合もあるためである。

したがって、配置値の比較と、配置集合そのものの比較は区別する。
Kを配置・関係・分離を統合した指標として完成させるには、
\[
    \mathcal{Z}_a,\qquad
    \mathcal{Z}_a\cap\mathcal{Z}_b,\qquad
    \mathcal{Z}_a\setminus\mathcal{Z}_b,\qquad
    \mathcal{Z}_a\triangle\mathcal{Z}_b
\]
によって得られる情報を、数値化の段階でどのように組み合わせるかを
明示する必要がある。

この統合を表す一般関数を導入する場合には、
\[
    K(a,b)
    =
    W_K\bigl(
        \mathcal{Z}_a,
        \mathcal{Z}_b,
        \mathcal{Z}_a\cap\mathcal{Z}_b,
        \mathcal{Z}_a\setminus\mathcal{Z}_b,
        \mathcal{Z}_b\setminus\mathcal{Z}_a
    \bigr)
\]
のような形式を考えることはできる。
ただし、これは関数の入力に何を含めるかを示すための形式的な表記であり、
$W_K$ の具体的な関数形を確定するものではない。
$W_K$ の具体形を採用する場合には、その定義域、値域、
および各集合演算が数値に反映される方法を明示しなければならない。

\subsection{対称差の測度による数値化}

$\mathcal{Z}_a$、$\mathcal{Z}_b$ が可測で、
$\mu(\mathcal{Z})$ が有限かつ正である場合には、
配置間の分離を対称差の正規化測度によって数値化できる。

\[
    d_\mu(a,b)
    =
    \frac{
        \mu(\mathcal{Z}_a\triangle\mathcal{Z}_b)
    }{
        \mu(\mathcal{Z})
    }
\]

このとき、
\[
    0\leq d_\mu(a,b)\leq1
\]
が成り立つ。

さらに、対称差の測度には
\[
\begin{aligned}
    \mu(\mathcal{Z}_a\triangle\mathcal{Z}_b)
    &=
    \mu(\mathcal{Z}_a)
    +\mu(\mathcal{Z}_b)\\
    &\quad
    -2\mu(\mathcal{Z}_a\cap\mathcal{Z}_b)
\end{aligned}
\]
という関係がある。
したがって、
\[
    d_\mu(a,b)
    =
    K(a)+K(b)
    -2\frac{\mu(\mathcal{Z}_a\cap\mathcal{Z}_b)}
             {\mu(\mathcal{Z})}
\]
となる。

この式は、個々の配置値だけでなく、
両者の共通配置を含めて分離を数値化するものである。
ただし、これをK全体の最終定義と同一視するのではなく、
配置間の分離を数値化する一つの構成要素として扱う。

また、$d_\mu(a,b)=0$ であっても、
必ずしも $\mathcal{Z}_a=\mathcal{Z}_b$ とは限らない。
この場合に保証されるのは、両者の対称差が測度ゼロであることまでである。

\subsection{方向別の差と対称差の区別}

配置の方向別の違いを数値化する場合には、
\[
    d_\mu^+(a\mid b)
    =
    \frac{\mu(\mathcal{Z}_a\setminus\mathcal{Z}_b)}
         {\mu(\mathcal{Z})},
\]
\[
    d_\mu^-(a\mid b)
    =
    \frac{\mu(\mathcal{Z}_b\setminus\mathcal{Z}_a)}
         {\mu(\mathcal{Z})}
\]
を考えることができる。
ここでも、各差集合が可測であり、
正規化の条件が満たされていることを仮定する。

このとき、
\[
    d_\mu(a,b)
    =
    d_\mu^+(a\mid b)+d_\mu^-(a\mid b)
\]
が成り立つ。

対称差の数値は、どちらの公理だけに配置されているかという方向性を
単独では保持しない。
そのため、方向別の差と対称差の数値は、
必要に応じて区別して扱う。

\subsection{区間評価における許容誤差の条件}

区間
\[
    K(a)\in[k_a^-,k_a^+]
\]
に対して、
\[
    \Delta K(a)=k_a^+-k_a^-
\]
と定義する。

事前に定めた許容誤差を $\varepsilon>0$ とするとき、
\[
    \Delta K(a)\leq\varepsilon
\]
を、単一の近似値を採用するための幅に関する条件とする。

ただし、この条件は、比較対象間の区別が保存されることを
単独で保証するものではない。
比較対象間の区別を保証するには、
\[
    \Delta_{\max}<d_{\min}^-(A_0)
\]
のような最小差に基づく条件を別途確認する。

したがって、許容誤差を満たすことと、
比較対象間の正の差が保証されることは異なる条件である。
前者を満たしても後者を満たさない場合には、
近似値を使用する際に区別可能性が保証されないことを明示する。

\subsection{数値化に関する結論}

Kの数値化では、次の区別を維持する。

\begin{enumerate}[label=(\arabic*)]
    \item 公理ごとの配置割合を表す数値。
    \item Solver内の複数の配置値から求める最小値・平均値・最大値。
    \item 配置集合間の共通部分・差集合・対称差に基づく関係と分離。
    \item 区間評価によって保証される値の範囲。
    \item 許容誤差の条件を満たす場合に限って採用する近似値。
    \item $B$、$I$ と組み合わせて得られる統合評価
    $T_{\mathrm{overlay}}$。
\end{enumerate}

これらを区別することにより、
集合論的な配置構造、数値として保証できる内容、
および近似によって得られる内容を混同せずに扱うことができる。

\end{document}